
\subsubsection{3.1 Research Questions}

\begin{itemize}
\item \textbf{RQ1:} Does pretraining paradigm significantly modulate the spurious/task probe accuracy ratio on frozen ResNet-50 features (Waterbirds)?
\item \textbf{RQ2:} Is this effect consistent across datasets with different spurious attribute types, and does the paradigm ranking change (CelebA)?
\item \textbf{RQ3:} Is augmentation a functional causal mechanism for modulating spurious feature encoding in contrastive SSL?
\end{itemize}

\subsubsection{3.2 Pretraining Paradigms}

Four ResNet-50 backbones pretrained on ImageNet-1k under different objectives:

\begin{itemize}
\item \textbf{ERM (Supervised):} Standard cross-entropy classification (\texttt{torchvision resnet50(pretrained=True)}). ImageNet top-1: 76.1\%.
\item \textbf{MoCo-v3 (Contrastive SSL):} Momentum-contrast contrastive SSL \citep{He2020MoCo}; official \texttt{r-50-1000ep.pth.tar} checkpoint loaded via custom \texttt{hubconf.py} (PyTorch Hub download was unavailable). ImageNet top-1: 74.3\%.
\item \textbf{DINO (Self-distillation):} Self-distillation with momentum teacher \citep{Caron2021DINO}; \texttt{dino\textbackslash \_resnet50} via PyTorch Hub. ImageNet top-1: 75.3\%.
\item \textbf{BarlowTwins (Non-contrastive SSL):} Redundancy-reduction via cross-correlation matrix alignment \citep{Zbontar2021BarlowTwins}; official weights via direct download. ImageNet top-1: 73.5\%.
\end{itemize}

All backbones output 2048-dimensional features via global average pooling (fc = Identity) and are frozen throughout all probe experiments.

\subsubsection{3.3 Spurious/Task Probe Accuracy Ratio}

For each frozen backbone, separate linear probes (logistic regression, C = 1.0, lbfgs solver, max\_iter = 1000) predict the spurious attribute and the task label from the 2048-dimensional features:

$$\text{ratio} = \frac{\text{spurious\_probe\_acc}}{\text{task\_probe\_acc}}$$

A ratio above 1.0 indicates that the spurious attribute is more linearly decodable than the task label. This scale-free metric decouples paradigm effects from absolute accuracy differences across datasets.

\subsubsection{3.4 Group-Balanced Probe Protocol}

Probes are trained on a group-balanced subset of the validation split (equal examples per task label $\times$ spurious attribute group), not the training split. This avoids inflating spurious probe accuracy due to the 95\% Waterbirds training correlation. Evaluation uses the full balanced test split (Waterbirds: 5,794 images; CelebA: 720 images). Five random seeds per probe (seeds 0--4).

\subsubsection{3.5 Statistical Analysis}

Pairwise comparisons use two-sample t-tests across five seeds, Bonferroni-corrected for six pairs (C(4,2)). Gate criterion: p\_bonf < 0.05 and mean difference $\geq$ 0.02. Cohen's d is computed with pooled standard deviation (d = ($\mu _1 - \mu _{2})/s_pooled$, s\_pooled = $\sqrt((s_{1}^2 + s_{2}^{2})/2))$. One-way ANOVA tests joint paradigm significance.

\subsubsection{3.6 Background-Replacement Augmentation Ablation}

SimCLR trained from scratch on Waterbirds under two conditions: \textbf{SimCLR-Original} (standard augmentation: random resized crop, horizontal flip, color jitter, Gaussian blur) and \textbf{SimCLR-NoBackground} (adds random background replacement using Places365 images drawn from a 10,000-image pool, applied to CUB-200-2011 segmented foregrounds). Both conditions share matched hyperparameters: SGD (momentum = 0.9, weight\_decay = 1e-4), learning rate = 0.03 (cosine annealing, T\_max = 50), batch size = 256, 50 epochs, NT-Xent temperature = 0.5, projection head 2048 $\rightarrow$ 2048 $\rightarrow$ 128 (L2-normalized). Five seeds per condition.

Mechanism activation is verified by measuring pixel difference between the background region of the two augmented views: if the background was replaced independently in each view, pixel difference in the background region should far exceed the threshold (0.05). Achieved: pixel\_diff = 0.9656 ($19\times$ threshold).

\subsubsection{3.7 Datasets}

\textbf{Waterbirds} \citep{Sagawa2020GroupDRO}: 4,795 training images. Balanced test: 5,794 images (50\% spurious per class). Task = bird species (two classes); spurious = background type (land/water; 95\% training correlation).

\textbf{CelebA} \citep{Liu2015CelebA}: 162,770 training images. Balanced test: 720 images (180 per group $\times$ 4 groups). Task = Blond\_Hair (attribute column 9); spurious = Male (attribute column 20; approximately 85\% training correlation).

